\documentclass[10pt,a4paper]{amsart}
\usepackage{amsmath,amssymb,amsthm,graphicx}
% Retained minimal preamble for the Stacks Project chapter more-algebra.tex (More on Algebra).
% The chapter file pinned-original-source.tex carries only \input{preamble} on its first line, so no
% preamble text lives inside the chapter itself.  The Stacks Project uses one single shared preamble for
% every chapter; the snapshot stages it identically as sources/stacks-chow/preamble.tex and
% sources/stacks-etale-cohomology/preamble.tex.  The block below is copied mechanically, line for line,
% from sources/stacks-chow/preamble.tex lines 14-18 (the Xy-pic setup) and 154-217 (the native theorem
% declarations together with the mathematical macros).  Nothing was added, removed or redefined; the
% wrapper class and the amsmath/amssymb/amsthm/graphicx base packages are supplied by the assembler.
% Chapter: more-algebra.tex -- More on Algebra.
% For commutative diagrams we use Xy-pic
\usepackage[all]{xy}

% We use 2cell for 2-commutative diagrams.
\xyoption{2cell}

% Theorem environments.
%
\theoremstyle{plain}
\newtheorem{theorem}[subsection]{Theorem}
\newtheorem{proposition}[subsection]{Proposition}
\newtheorem{lemma}[subsection]{Lemma}

\theoremstyle{definition}
\newtheorem{definition}[subsection]{Definition}
\newtheorem{example}[subsection]{Example}
\newtheorem{exercise}[subsection]{Exercise}
\newtheorem{situation}[subsection]{Situation}

\theoremstyle{remark}
\newtheorem{remark}[subsection]{Remark}
\newtheorem{remarks}[subsection]{Remarks}

\numberwithin{equation}{subsection}

% Macros
%
\def\lim{\mathop{\mathrm{lim}}\nolimits}
\def\colim{\mathop{\mathrm{colim}}\nolimits}
\def\Spec{\mathop{\mathrm{Spec}}}
\def\Hom{\mathop{\mathrm{Hom}}\nolimits}
\def\Ext{\mathop{\mathrm{Ext}}\nolimits}
\def\SheafHom{\mathop{\mathcal{H}\!\mathit{om}}\nolimits}
\def\SheafExt{\mathop{\mathcal{E}\!\mathit{xt}}\nolimits}
\def\Sch{\mathit{Sch}}
\def\Mor{\mathop{\mathrm{Mor}}\nolimits}
\def\Ob{\mathop{\mathrm{Ob}}\nolimits}
\def\Sh{\mathop{\mathit{Sh}}\nolimits}
\def\NL{\mathop{N\!L}\nolimits}
\def\CH{\mathop{\mathrm{CH}}\nolimits}
\def\proetale{{pro\text{-}\acute{e}tale}}
\def\etale{{\acute{e}tale}}
\def\QCoh{\mathit{QCoh}}
\def\Ker{\mathop{\mathrm{Ker}}}
\def\Im{\mathop{\mathrm{Im}}}
\def\Coker{\mathop{\mathrm{Coker}}}
\def\Coim{\mathop{\mathrm{Coim}}}

% Boxtimes
%
\DeclareMathSymbol{\boxtimes}{\mathbin}{AMSa}{"02}

%
% Macros for moduli stacks/spaces
%
\def\QCohstack{\mathcal{QC}\!\mathit{oh}}
\def\Cohstack{\mathcal{C}\!\mathit{oh}}
\def\Spacesstack{\mathcal{S}\!\mathit{paces}}
\def\Quotfunctor{\mathrm{Quot}}
\def\Hilbfunctor{\mathrm{Hilb}}
\def\Curvesstack{\mathcal{C}\!\mathit{urves}}
\def\Polarizedstack{\mathcal{P}\!\mathit{olarized}}
\def\Complexesstack{\mathcal{C}\!\mathit{omplexes}}
% \Pic is the operator that assigns to X its picard group, usage \Pic(X)
% \Picardstack_{X/B} denotes the Picard stack of X over B
% \Picardfunctor_{X/B} denotes the Picard functor of X over B
\def\Pic{\mathop{\mathrm{Pic}}\nolimits}
\def\Picardstack{\mathcal{P}\!\mathit{ic}}
\def\Picardfunctor{\mathrm{Pic}}
\def\Deformationcategory{\mathcal{D}\!\mathit{ef}}

% --- appended 2026-10-12 (Chief): the project own definitions of the reference/slogan/history
% display environments, copied mechanically from the shared Stacks preamble lines 8-12
% (\usepackage{verbatim} plus the three \newenvironment lines, which the chapter preamble keeps
% verbatim from the project preamble). Nothing added, removed or redefined: these environments
% carry no mathematical content and the project itself typesets them as comments. Without these
% three lines a counted statement whose own text uses them fails with Environment slogan undefined.
\usepackage{verbatim}
\newenvironment{reference}{\comment}{\endcomment}
\newenvironment{slogan}{\comment}{\endcomment}
\newenvironment{history}{\comment}{\endcomment}
\usepackage{cleveref}
\crefname{originallemma}{original lemma}{original lemmas}
\Crefname{originallemma}{Original lemma}{Original lemmas}
\makeatletter
\newlabel{more-algebra-lemma-diagonal-better}{{0BYM}{source}{}{sourcepointer-1}{}}
\newlabel{more-algebra-lemma-diagonal-better@cref}{{[originallemma][0][]0BYM}{[1][0][]source}{}{}{}}
\makeatother
\crefname{originallemma}{original lemma}{original lemmas}
\Crefname{originallemma}{Original lemma}{Original lemmas}
\makeatletter
\newlabel{remark-internal-hom}{{0GWQ}{source}{}{sourcepointer-2}{}}
\newlabel{remark-internal-hom@cref}{{[originallemma][0][]0GWQ}{[1][0][]source}{}{}{}}
\makeatother
\crefname{originallemma}{original lemma}{original lemmas}
\Crefname{originallemma}{Original lemma}{Original lemmas}
\makeatletter
\newlabel{categories-remark-internal-hom-monoidal}{{0GWJ}{source}{}{sourcepointer-3}{}}
\newlabel{categories-remark-internal-hom-monoidal@cref}{{[originallemma][0][]0GWJ}{[1][0][]source}{}{}{}}
\makeatother
\title{A canonical morphism from a total tensor complex into an internal hom complex}
\author{The Stacks Project Authors}
\date{stacks-project:more-algebra.tex:local-snapshot-20261012; GNU FDL 1.2 or later}
\begin{document}
\maketitle
\setcounter{subsection}{1}
\noindent\textit{Editorial note.} All delivered mathematical text is the author's own, copied from the pinned Stacks Project chapter file more-algebra.tex; no statement or proof step was rewritten, reordered or removed.\par\medskip
\noindent\textit{Editorial note.} The chapter file pinned-original-source.tex carries only \textbackslash{}input\{preamble\} on its first line, so no preamble text lives inside the chapter. The retained file sources/stacks-more-algebra/retained-preamble.tex is a mechanical, line-for-line copy of the project's single shared preamble as staged in the snapshot (Xy-pic setup plus theorem declarations and mathematical macros). Nothing was added, removed or redefined, and the excerpt serves only as preamble, never as a span coordinate.\par\medskip
\noindent\textit{Editorial note.} Both counted spans start at the \textbackslash{}begin line of their environment, so the statement and the author's proof are delivered whole; surrounding prose of the section is deliberately outside the counted spans and is carried as standard prerequisites instead.\par\medskip
\noindent\textit{Editorial note.} The citation fields (source\_id, source\_version, authors, source\_url, license.name, source\_title) are identical to those of every other card this lane has produced from this chapter, so the batch selector treats them as one source.\par\medskip
\noindent\textit{Allowed standard prerequisites (editorial).}\begin{itemize}
\item Internal hom complexes and their compatibility with the monoidal structure. Remark \textbackslash{}ref\{remark-internal-hom\} and Categories, Remark \textbackslash{}ref\{categories-remark-internal-hom-monoidal\}, cited at proof lines 18048-18050. Supplies the canonical map abstractly; the author then constructs it directly.
\item Sign rules for the differential of the total tensor complex and of internal hom complexes. used at proof lines 18059-18067, 18083-18091, and 18099-18105. Give the two differential computations whose compatibility the construction must satisfy.
\end{itemize}\medskip
\ifdefined\hypertarget\hypertarget{sourcepointer-1}{}\fi
\noindent\textit{Source reference.} Original lemma 0BYM: focused-seeds/AGP/sources/stacks-snapshot-20261011/tags.txt, source lines 15020--15020. The lemma's own canonical tag, quoted so that readers can resolve the label inside the delivered text; it points at the counted statement itself and carries no proof content.\par\medskip
\ifdefined\hypertarget\hypertarget{sourcepointer-2}{}\fi
\noindent\textit{Source reference.} Original lemma 0GWQ: focused-seeds/AGP/sources/stacks-snapshot-20261011/tags.txt, source lines 21080--21080. Navigation pointer to a statement the counted proof refers to but does not contain; nothing load-bearing is delegated to it.\par\medskip
\ifdefined\hypertarget\hypertarget{sourcepointer-3}{}\fi
\noindent\textit{Source reference.} Original lemma 0GWJ: focused-seeds/AGP/sources/stacks-snapshot-20261011/tags.txt, source lines 21074--21074. Navigation pointer to a statement the counted proof refers to but does not contain; nothing load-bearing is delegated to it.\par\medskip
\begin{lemma}
\label{lemma-diagonal-better}
Let $R$ be a ring. Given complexes $K^\bullet, L^\bullet, M^\bullet$
of $R$-modules there is a canonical morphism
$$
\text{Tot}(K^\bullet \otimes_R \Hom^\bullet(M^\bullet, L^\bullet))
\longrightarrow
\Hom^\bullet(M^\bullet, \text{Tot}(K^\bullet \otimes_R L^\bullet))
$$
of complexes of $R$-modules functorial in all three complexes.
\end{lemma}

\begin{proof}
Via the discussion in Remark \ref{remark-internal-hom} the existence
of such a canonical map follows from
Categories, Remark \ref{categories-remark-internal-hom-monoidal}.
We also give a direct construction.

\medskip\noindent
Let $\alpha$ be an element of degree $n$ of the right hand side.
Thus
$$
\alpha = (\alpha^{p, q}) \in \prod\nolimits_{p + q = n}
\Hom_R(M^{-q}, \text{Tot}^p(K^\bullet \otimes_R L^\bullet))
$$
Each $\alpha^{p, q}$ is an element
$$
\alpha^{p, q} = (\alpha^{r, s, q}) \in
\Hom_R(M^{-q}, \bigoplus\nolimits_{r + s + q = n} K^r \otimes_R L^s)
$$
where we think of $\alpha^{r, s, q}$ as a family of maps such that
for every $x \in M^{-q}$ only a finite number of
$\alpha^{r, s, q}(x)$ are nonzero. By our sign rules we get
\begin{align*}
\text{d}(\alpha^{r, s, q})
& =
\text{d}_{\text{Tot}(K^\bullet \otimes_R L^\bullet)} \circ \alpha^{r, s, q}
- (-1)^n \alpha^{r, s, q} \circ \text{d}_M \\
& =
\text{d}_K \circ \alpha^{r, s, q} + (-1)^r \text{d}_L \circ \alpha^{r, s, q}
- (-1)^n \alpha^{r, s, q} \circ \text{d}_M
\end{align*}
On the other hand, if $\beta$ is an element of degree $n$ of the
left hand side, then
$$
\beta = (\beta^{p, q}) \in
\bigoplus\nolimits_{p + q = n} K^p \otimes_R \Hom^q(M^\bullet, L^\bullet)
$$
and we can write $\beta^{p, q} = \sum \gamma_i^p \otimes \delta_i^q$ with
$\gamma_i^p \in K^p$ and
$$
\delta_i^q = (\delta_i^{r, s}) \in
\prod\nolimits_{r + s = q} \Hom_R(M^{-s}, L^r)
$$
By our sign rules we have
\begin{align*}
\text{d}(\beta^{p, q})
& =
\text{d}_K(\beta^{p, q}) +
(-1)^p \text{d}_{\Hom^\bullet(M^\bullet, L^\bullet)}(\beta^{p, q}) \\
& =
\sum \text{d}_K(\gamma_i^p) \otimes \delta_i^q +
(-1)^p \sum \gamma_i^p \otimes
(\text{d}_L \circ \delta_i^q - (-1)^q \delta_i^q \circ \text{d}_M)
\end{align*}
We send the element $\beta$ to $\alpha$ with
$$
\alpha^{r, s, q} = c^{r, s, q}(\sum \gamma_i^r \otimes \delta_i^{s, q})
$$
where $c^{r, s, q} : K^r \otimes_R \Hom_R(M^{-q}, L^s) \to
\Hom_R(M^{-q}, K^r \otimes_R L^s)$ is the canonical map.
For a given $\beta$ and $r$ there are only finitely many nonzero
$\gamma_i^r$ hence only finitely many nonzero $\alpha^{r, s, q}$
are nonzero (for a given $r$). Thus this family of maps satisfies
the conditions above and the map is well defined.
Comparing signs we see that this is compatible with differentials.
\end{proof}

\end{document}
